\documentclass[11pt,a4paper]{article}

% ---------------------------------------------------------
% PREAMBLE (stable, warning-free, Overleaf-safe)
% ---------------------------------------------------------
\usepackage[utf8]{inputenc}
\usepackage[T1]{fontenc}
\usepackage{lmodern}
\usepackage{amsmath, amssymb, amsfonts}
\usepackage{bm}
\usepackage{geometry}
\usepackage{graphicx}
\usepackage{hyperref}
\usepackage{cite}
\usepackage{authblk}
\usepackage{physics}
\usepackage{mathtools}

\geometry{margin=2.4cm}

\title{\textbf{Companion LXXX: Neural Schrödinger Flows and Score-Transport Geometry of Persistence Diagrams in the Relaxation-Driven Cyclic Cosmology}}
\author{Michael Lehmann \\ \texttt{mi.lehmann@gmx.de}}
\date{April 2026}

\begin{document}
\maketitle

\begin{abstract}
Neural Schrödinger flows unify score-based diffusion models and Schrödinger bridge 
dynamics into a single framework for learning stochastic transport processes. 
Applied to persistence diagrams, these flows learn the time-dependent drift and 
diffusion fields that interpolate between topological ensembles in a probabilistic 
and geometrically meaningful way. Building on the entropic optimal transport and 
score-based models developed in previous Companions, this work introduces neural 
Schrödinger flows for persistence diagrams in the Relaxation-Driven Cyclic 
Cosmology (RDCC). The relaxation mechanism modifies the scale-dependent evolution 
of the gravitational potential, inducing characteristic changes in score fields, 
transport geometry, and stochastic flow trajectories. The resulting framework 
provides a powerful and flexible tool for modeling RDCC signatures in weak-lensing 
topology.
\end{abstract}
% ---------------------------------------------------------
\section{Introduction}

Persistence diagrams encode the birth and death of topological features in weak-
lensing convergence fields. Previous Companions introduced Wasserstein geometry, 
entropic optimal transport, Schrödinger bridges, and score-based diffusion models 
as tools for analyzing the Relaxation-Driven Cyclic Cosmology (RDCC).

Neural Schrödinger flows unify these approaches by learning the drift and 
diffusion fields of a stochastic differential equation (SDE) that transports one 
distribution of persistence diagrams into another. This combines the dynamic 
interpretation of Schrödinger bridges with the generative power of score-based 
models, yielding a flexible and expressive framework for modeling topological 
distributions.

In RDCC, the relaxation mechanism modifies the underlying gravitational potential 
in a scale-dependent manner, producing characteristic signatures in the learned 
transport geometry.
% ---------------------------------------------------------
\section{Neural Schrödinger Flows}

A neural Schrödinger flow is defined by the SDE

\begin{equation}
    dD_{t}
    = v_{\theta}(D_{t}, t) \, dt
    + \sqrt{2 \varepsilon(t)} \, dB_{t},
\end{equation}

where

\begin{itemize}
    \item $v_{\theta}(D_{t}, t)$ is a neural network approximating the drift,
    \item $\varepsilon(t)$ is a time-dependent diffusion schedule,
    \item $B_{t}$ is Brownian motion on diagram space.
\end{itemize}

The flow is trained to match boundary distributions $\rho_{0}$ and $\rho_{1}$ by 
minimizing a Schrödinger-bridge-inspired objective.
% ---------------------------------------------------------
\section{Score-Transport Geometry}

The drift is decomposed as

\begin{equation}
    v_{\theta}(D, t)
    = - \varepsilon(t) \, s_{\theta}(D, t)
      + u_{\theta}(D, t),
\end{equation}

where

\begin{itemize}
    \item $s_{\theta}(D, t)$ is the learned score field,
    \item $u_{\theta}(D, t)$ is a learned residual transport term.
\end{itemize}

This decomposition yields a geometric interpretation:

\begin{itemize}
    \item the score term drives the flow toward high-density regions,
    \item the residual term encodes transport geometry beyond pure diffusion.
\end{itemize}
% ---------------------------------------------------------
\section{RDCC Modification of Neural Schrödinger Dynamics}

In RDCC, the convergence evolves as

\begin{equation}
    \kappa_{\mathrm{RDCC}}(\ell)
    = \kappa(\ell)
      \left[
        1 - \chi_{\kappa}
        \left(\frac{\ell}{\ell_{\mathrm{rel}}}\right)^{\alpha}
      \right].
\end{equation}

This modifies the log-density of persistence diagrams:

\begin{equation}
    \log p_{\mathrm{RDCC}}(D)
    = \log p(D)
      - \chi_{p}
        \left(\frac{\ell}{\ell_{\mathrm{rel}}}\right)^{\alpha}.
\end{equation}

The RDCC score becomes

\begin{equation}
    s_{\mathrm{RDCC}}(D, t)
    = s(D, t)
      \left[
        1 - \chi_{s}
        \left(\frac{\ell}{\ell_{\mathrm{rel}}}\right)^{\alpha}
      \right].
\end{equation}

The RDCC drift is

\begin{equation}
    v^{\mathrm{RDCC}}_{\theta}(D, t)
    = v_{\theta}(D, t)
      \left[
        1 - \chi_{v}
        \left(\frac{\ell}{\ell_{\mathrm{rel}}}\right)^{\alpha}
      \right].
\end{equation}
% ---------------------------------------------------------
\section{RDCC Signatures in Neural Transport Geometry}

RDCC induces:

\begin{itemize}
    \item reduced score magnitude at small scales,
    \item enhanced drift at large scales,
    \item modified stochastic flow trajectories,
    \item altered dynamic barycenters,
    \item scale-dependent deformation of transport geometry,
    \item characteristic curvature in learned Schrödinger flows.
\end{itemize}

These signatures provide a generative and dynamic probe of the RDCC gravitational 
field.
% ---------------------------------------------------------
\section{Conclusion}

Neural Schrödinger flows unify score-based diffusion models and Schrödinger bridge 
dynamics into a single framework for learning stochastic transport processes. In 
the Relaxation-Driven Cyclic Cosmology, the relaxation mechanism modifies the 
score field, drift dynamics, and transport geometry in a scale-dependent manner, 
producing distinctive signatures in weak-lensing topology. The present Companion 
establishes the foundation for future studies of neural SDEs, generative modeling, 
and dynamic geometric analysis in cosmological topology.

\bibliographystyle{apsrev4-2}
\nocite{*}
\bibliography{references}


\end{document}
